\documentclass[11pt]{article}
\usepackage[margin=24mm]{geometry}\usepackage[T1]{fontenc}
\usepackage{amsmath,amssymb,amsthm,mathtools,mathrsfs,longtable,booktabs,array,graphicx,xurl}
\usepackage[hidelinks,hypertexnames=false]{hyperref}
\allowdisplaybreaks\setlength{\emergencystretch}{3em}
\providecommand{\tightlist}{\setlength{\itemsep}{0pt}\setlength{\parskip}{0pt}}
\DeclareUnicodeCharacter{03C0}{\ensuremath{\pi}}
\DeclareUnicodeCharacter{03C7}{\ensuremath{\chi}}
\DeclareUnicodeCharacter{03A6}{\ensuremath{\Phi}}
\DeclareUnicodeCharacter{0394}{\ensuremath{\Delta}}
\DeclareUnicodeCharacter{0393}{\ensuremath{\Gamma}}
\DeclareUnicodeCharacter{039B}{\ensuremath{\Lambda}}
\DeclareUnicodeCharacter{03C3}{\ensuremath{\sigma}}
\DeclareUnicodeCharacter{03B7}{\ensuremath{\eta}}
\DeclareUnicodeCharacter{03C4}{\ensuremath{\tau}}
\DeclareUnicodeCharacter{03B4}{\ensuremath{\delta}}
\DeclareUnicodeCharacter{03B3}{\ensuremath{\gamma}}
\DeclareUnicodeCharacter{03B1}{\ensuremath{\alpha}}
\DeclareUnicodeCharacter{03B2}{\ensuremath{\beta}}
\DeclareUnicodeCharacter{03C1}{\ensuremath{\rho}}
\DeclareUnicodeCharacter{03BA}{\ensuremath{\kappa}}
\DeclareUnicodeCharacter{03BB}{\ensuremath{\lambda}}
\DeclareUnicodeCharacter{03B8}{\ensuremath{\theta}}
\DeclareUnicodeCharacter{03B5}{\ensuremath{\varepsilon}}
\DeclareUnicodeCharacter{2208}{\ensuremath{\in}}
\DeclareUnicodeCharacter{2264}{\ensuremath{\le}}
\DeclareUnicodeCharacter{2265}{\ensuremath{\ge}}
\DeclareUnicodeCharacter{2260}{\ensuremath{\ne}}
\DeclareUnicodeCharacter{2192}{\ensuremath{\to}}
\DeclareUnicodeCharacter{2282}{\ensuremath{\subset}}
\DeclareUnicodeCharacter{2286}{\ensuremath{\subseteq}}
\DeclareUnicodeCharacter{2297}{\ensuremath{\otimes}}
\DeclareUnicodeCharacter{03A3}{\ensuremath{\Sigma}}
\DeclareUnicodeCharacter{221A}{\ensuremath{\surd}}
\DeclareUnicodeCharacter{221E}{\ensuremath{\infty}}
\DeclareUnicodeCharacter{00B2}{\ensuremath{{}^2}}
\DeclareUnicodeCharacter{00B3}{\ensuremath{{}^3}}
\DeclareUnicodeCharacter{2074}{\ensuremath{{}^4}}
\DeclareUnicodeCharacter{2080}{\ensuremath{{}_0}}
\DeclareUnicodeCharacter{2081}{\ensuremath{{}_1}}
\DeclareUnicodeCharacter{2082}{\ensuremath{{}_2}}
\DeclareUnicodeCharacter{2096}{\ensuremath{{}_k}}
\DeclareUnicodeCharacter{2212}{\ensuremath{-}}
\DeclareUnicodeCharacter{00D7}{\ensuremath{\times}}
\DeclareUnicodeCharacter{00B1}{\ensuremath{\pm}}
\DeclareUnicodeCharacter{2202}{\ensuremath{\partial}}
\DeclareUnicodeCharacter{2225}{\ensuremath{\Vert}}
\DeclareUnicodeCharacter{2113}{\ensuremath{\ell}}
\DeclareUnicodeCharacter{211D}{\ensuremath{\mathbb R}}
\DeclareUnicodeCharacter{2102}{\ensuremath{\mathbb C}}
\DeclareUnicodeCharacter{03BC}{\ensuremath{\mu}}
\DeclareUnicodeCharacter{03B6}{\ensuremath{\zeta}}
\DeclareUnicodeCharacter{03BE}{\ensuremath{\xi}}
\DeclareUnicodeCharacter{220F}{\ensuremath{\prod}}
\begin{document}
\section{Independent derivation of the complete Gamma reduction for the original kernel}\label{independent-derivation-of-the-complete-gamma-reduction-for-the-original-kernel}

Date: 2026-09-20. Checked source: \path{JOINT_MINIMUM_RECEIVERS.tex}, JM21--JM28, in this directory. All constants and all displayed inequalities in those equations are accepted. The only textual correction is the reference to P5: P5 is the tail estimate, whereas the multiplier estimate used here is the Gamma multiplication estimate in its proof, written explicitly as C4 in \path{COERCIVITY_INDEPENDENT.md}.

The result retains the actual fixed kernel of the original observation. Its dimension \texttt{m\_K=8k-16} and the simple-packet formula \texttt{q=(k+1)\^{}2} are used on their original five-orbit domain. No kernel frame is selected or altered in this derivation. The asymptotic limit is along that domain with \texttt{k\ congruent\ to\ 1\ modulo\ 4}.

\subsection{1. Exact convolution density and its mass}\label{exact-convolution-density-and-its-mass}

Keep \[
\sigma(y)=\frac{|\Gamma(1/4+iy/2)|^2}{2\pi},\quad
w_h(y)=\frac{|(2\xi/h)(1/2+iy)|^2}{2\pi},\quad
M_\sigma=\sqrt{2\pi}.
\] Neither measure is divided by its mass. To check every convolution constant, use the Fourier convention \[
\widehat f(t)=\int_{\mathbb R}e^{-ity}f(y)\,dy.
\] For any a\textgreater0 the beta integral gives \[
\frac{\Gamma(a+ix)\Gamma(a-ix)}{\Gamma(2a)}
=\int_{-\infty}^{\infty}e^{ixv}(2\cosh(v/2))^{-2a}\,dv.
\tag{KG1}
\] For completeness, multiply the two Euler integrals for Gamma and set \texttt{r=u+v}, \texttt{theta=u/(u+v)} in the positive quadrant. The Jacobian is r; the r integral is Gamma(2a), and the remaining integral is \texttt{integral\_0\^{}1\ theta\^{}\{a+ix-1\}(1-theta)\^{}\{a-ix-1\}\ dtheta}. Then set \texttt{z=log(theta/(1-theta))}, whose derivative is \texttt{dtheta/dz=theta(1-theta)}. The resulting integrand is precisely the right side of KG1. Absolute convergence is ensured by a\textgreater0.

The function \texttt{(2cosh(v/2))\^{}\{-2a\}} and its first two derivatives are integrable. Its Fourier transform is therefore integrable: on a bounded interval it is bounded by the L1 norm, and outside that interval two integrations by parts give a constant times \textbar x\textbar\^{}\{-2\}. Fourier inversion thus applies to KG1 as an ordinary absolutely convergent integral. After the exact substitution y=2x it gives \[
\int_{\mathbb R}e^{-ity}
\frac{|\Gamma(a+iy/2)|^2}{2\pi}\,dy
=\frac{2\Gamma(2a)}{(2\cosh t)^{2a}}.
\tag{KG2}
\] In particular \[
\widehat\sigma(t)=M_\sigma(\cosh t)^{-1/2},\qquad
\int\sigma=M_\sigma.
\] For the k-fold convolution, multiplication of Fourier transforms and KG2 with a=k/4 give \[
\sigma^{*k}(y)=
\frac{M_\sigma^k2^{k/2-1}}{\Gamma(k/2)}
\frac{|\Gamma(k/4+iy/2)|^2}{2\pi}.
\tag{KG3}
\] Both sides are integrable, and their Fourier transforms agree; uniqueness of the Fourier transform proves the equality. Integrating either side shows that the mass is exactly \texttt{M\_sigma\^{}k=(2pi)\^{}\{k/2\}}.

Now put \texttt{ell=(k-1)/4}. The Gamma recurrence gives \[
|\Gamma(1/4+\ell+iy/2)|^2
=4^{-\ell}\prod_{j=0}^{\ell-1}
 \left[y^2+(2j+1/2)^2\right]|\Gamma(1/4+iy/2)|^2.
\] Since \texttt{2\^{}\{k/2-1\}4\^{}\{-ell\}=2\^{}\{-1/2\}}, KG3 becomes \[
\boxed{\sigma^{*k}(y)=c_k\prod_{j<\ell}(y^2+b_j^2)\sigma(y),
\quad b_j=2j+\tfrac12,\quad
c_k=\frac{(2\pi)^{k/2}}{\sqrt2\,\Gamma(k/2)}.}
\tag{KG4}
\] This proves the exact JM21 constant. When k=1, ell=0 and c\_1=1; empty products therefore reproduce the original measure.

\subsection{2. The whole polynomial comparison}\label{the-whole-polynomial-comparison}

Set \texttt{n=2q}, \texttt{a=a\_\{h,n\}}, and \texttt{b=b\_\{h,n\}}, with exactly the constants proved in P1--P2 and C16--C19 of \path{COERCIVITY_INDEPENDENT.md}. The finite positive form inequality on one-factor degree at most n is \[
a\|Q\|_\sigma^2\le\|Q\|_{w_h}^2\le b\|Q\|_\sigma^2.
\] Tensor products of positive matrices preserve this order, with factors a\^{}k and b\^{}k. For a scalar polynomial P of degree at most n, retain the physical total variable \[
S=s_1+\cdots+s_k=k/2+i(y_1+\cdots+y_k)=c+iy,
\qquad c=k/2.
\] The polynomial \texttt{P(s\_1+...+s\_k)} has separate degree at most n. Its tensor norm is exactly its scalar convolution norm. Thus \[
a^k\|P(c+i\,\cdot)\|_{\sigma^{*k}}^2
\le\|P(c+i\,\cdot)\|_{w_h^{*k}}^2
\le b^k\|P(c+i\,\cdot)\|_{\sigma^{*k}}^2.
\tag{KG5}
\] This retains the complete form, including cross terms between every pair of coefficients and every relation candidate.

The monic Gamma basis has norms and recurrence \[
\gamma_j=M_\sigma j!(1/2)_j,\qquad
p_0=1,\quad p_1=y,\quad
p_{j+1}=yp_j-j(j-1/2)p_{j-1}.
\tag{KG6}
\] These formulas were checked from DLMF 18.19.8--9 and 18.23.7 in C1--C3 of \path{COERCIVITY_INDEPENDENT.md}, with the original measure retained. The two off-diagonal coefficients in its orthonormal basis are \texttt{sqrt(j(j-1/2))}. Consequently a polynomial Q of degree at most m satisfies \[
\|yQ\|_\sigma\le2(m+1)\|Q\|_\sigma.
\tag{KG7}
\] Indeed each of the two shifted coefficient vectors has norm at most \texttt{sqrt((m+1)(m+1/2))\textbar{}\textbar{}Q\textbar{}\textbar{}} or \texttt{sqrt(m(m-1/2))\textbar{}\textbar{}Q\textbar{}\textbar{}}; their sum is at most \texttt{2(m+1)\textbar{}\textbar{}Q\textbar{}\textbar{}}.

For \texttt{u\_j=2(n+j+1)+b\_j}, KG7 and the triangle inequality give \[
\|(y+ib_j)Q\|_\sigma\le u_j\|Q\|_\sigma
\quad\text{when }\deg Q\le n+j.
\] After j preceding factors, the degree is at most n+j. Iteration therefore gives exactly \[
\left\|\prod_{j<\ell}(y+ib_j)P(c+iy)\right\|_\sigma^2
\le\left(\prod_{j<\ell}u_j^2\right)
\|P(c+i\,\cdot)\|_\sigma^2.
\tag{KG8}
\] The matching lower bound uses the pointwise inequality \texttt{prod(y\^{}2+b\_j\^{}2)\textgreater{}=prod\ b\_j\^{}2}. Combining KG4, KG5, and KG8 proves \[
\boxed{\alpha_k\|P(c+i\,\cdot)\|_\sigma^2
\le\|P(c+i\,\cdot)\|_{w_h^{*k}}^2
\le\beta_k\|P(c+i\,\cdot)\|_\sigma^2,}
\tag{KG9}
\] \[
\alpha_k=a^kc_k\prod_{j<\ell}b_j^2,\qquad
\beta_k=b^kc_k\prod_{j<\ell}u_j^2.
\] This is JM22. The source measure on the right and left of KG9 is the order-one Gamma measure in the variable y, while the physical scalar coordinate remains S=c+iy throughout.

The finite logarithmic width is therefore \[
L_k^{\rm form}=\log(\beta_k/\alpha_k)
=k\log(b/a)+2\sum_{j<\ell}\log(u_j/b_j)\ge0.
\tag{KG10}
\] The scalar c\_k cancels in this ratio; it was included in both bounds before the cancellation. There is also an explicit upper bound for the second term. Since \[
\frac{u_j}{b_j}=2+\frac{2n+3/2}{2j+1/2}\le4n+5,
\] one has \[
0\le L_k^{\rm form}\le k\log(b/a)+2\ell\log(4n+5)
=O_h(k\log^2(q+2)).
\tag{KG11}
\] This verifies JM23, with no hidden scalar mass and no root-gap estimate.

\subsection{3. Quotient first, actual kernel restriction second}\label{quotient-first-actual-kernel-restriction-second}

Let E=C{[}S{]}/(chi\_k(S)) be the unchanged original scalar quotient of dimension q. For every \texttt{q-1\textless{}=N\textless{}=2q}, let \texttt{J\_N:P\_\{\textless{}=N\}-\textgreater{}E} be the original remainder map. For each v in E its affine fibre is nonempty. The complete physical minimum and Gamma minimum are \[
v^*G_Nv=\min_{J_NP=v}\|P(c+i\,\cdot)\|_{w_h^{*k}}^2,
\qquad
v^*G_N^\sigma v=\min_{J_NP=v}\|P(c+i\,\cdot)\|_\sigma^2.
\tag{KG12}
\] These minima exist because the spaces are finite-dimensional and their norm matrices are positive definite. Apply KG9 to every polynomial in the identical fibre. Taking the minima of the three forms gives \[
\alpha_kG_N^\sigma\preceq G_N\preceq\beta_kG_N^\sigma.
\tag{KG13}
\] No relation vector has been omitted, and no polynomial representative has been fixed before minimizing.

Let the original fixed kernel be \texttt{K=ker\ Lambda\_O}, and let \texttt{I\_K:C\^{}\{m\_K\}-\textgreater{}E} be its actual injective frame map. Set \[
H_{K,N}=I_K^*G_NI_K,\qquad
H_{K,N}^\sigma=I_K^*G_N^\sigma I_K.
\] Restriction of KG13 is a congruence by that same I\_K. Both restricted matrices are positive definite. The eigenvalues of \texttt{(H\_\{K,N\}\^{}sigma)\^{}\{-1/2\}H\_\{K,N\}(H\_\{K,N\}\^{}sigma)\^{}\{-1/2\}} lie in {[}alpha\_k,beta\_k{]}, so the product of its m\_K eigenvalues yields \[
m_K\log\alpha_k\le e_N\le m_K\log\beta_k,\qquad
e_N=\log\det H_{K,N}-\log\det H_{K,N}^\sigma.
\tag{KG14}
\] This is JM24. It is the restriction of the full quotient metric; it is not obtained by retaining selected coefficient directions before the quotient minimum.

Use the four original signs \[
\mathcal Re=e_{q-1}+e_q-e_{2q-1}-e_{2q}.
\] All four e\_N belong to the same interval of length m\_K L\_k\^{}\{form\}. Pair the first term with the third and the second with the fourth. Each difference has absolute value at most that length; hence \[
\boxed{\left|\mathcal R\log\det H_{K,N}
-\mathcal R\log\det H_{K,N}^\sigma\right|
\le2m_KL_k^{\rm form}.}
\tag{KG15}
\] No monotonicity of the errors e\_N is assumed. This verifies the factor 2 in JM25. The common scalar \texttt{m\_K\ log\ c\_k} in each e\_N cancels in this four-sign expression as well.

On the actual simple-packet domain, \[
2m_KL_k^{\rm form}=O_h(k^2\log^2(q+2)),\quad
\frac{k^2\log^2(q+2)}{kq}
=\frac{k\log^2((k+1)^2+2)}{(k+1)^2}\longrightarrow0.
\tag{KG16}
\] Thus the error is o\_h(kq), as stated, while retaining all roots and the entire original kernel.

\subsection{4. Exact covariance in the transported quotient}\label{exact-covariance-in-the-transported-quotient}

To specify the coordinate change completely, if chi\_k is monic of degree q, set \[
Q(y)=i^{-q}\chi_k(c+iy).
\] This is monic and retains every original root with its full multiplicity; the root lambda is transported to \texttt{(lambda-c)/i}. The substitution \[
T:E\longrightarrow E_y=\mathbb C[y]/(Q),\qquad
T[P]=[P(c+iy)]
\tag{KG17}
\] is an isomorphism, with inverse substitution \texttt{y=(S-c)/i}. Its coefficient matrix is triangular with diagonal \texttt{1,i,...,i\^{}\{q-1\}}. The transported kernel frame is \texttt{I\_\{K,y\}=T\ I\_K}. The matrices transform as \[
G_{N,y}^\sigma=(T^{-1})^*G_N^\sigma T^{-1},\qquad
C_{N,y}^\sigma=T C_N^\sigma T^*.
\] Consequently \texttt{I\_\{K,y\}\^{}*G\_\{N,y\}\^{}sigma\ I\_\{K,y\}=I\_K\^{}*G\_N\^{}sigma\ I\_K} exactly. No coordinate determinant or kernel-frame factor is lost.

In E\_y use the standard coefficient frame of degree below q and define \[
v_j=\operatorname{rem}_Qp_j,
\] where p\_j and gamma\_j are exactly KG6. The orthonormal source basis is \texttt{p\_j/sqrt(gamma\_j)}, for all \texttt{0\textless{}=j\textless{}=N}; its remainder matrix has columns \texttt{v\_j/sqrt(gamma\_j)}. Thus its complete covariance is \[
C_{N,y}^\sigma=\sum_{j=0}^N\frac{v_jv_j^*}{\gamma_j}.
\tag{KG18}
\] For j\textless q, p\_j already has degree below q, so v\_j=p\_j. The first q columns are triangular with diagonal one. They have full rank, proving \texttt{C\_\{N,y\}\^{}sigma\textgreater{}0} for every N\textgreater=q-1.

For a direct proof of the minimum formula, let B be this remainder matrix on orthonormal source coefficients. Since BB\emph{=C\textgreater0, the coefficient vector \texttt{B\^{}*C\^{}\{-1\}v} has value v and norm squared \texttt{v\^{}*C\^{}\{-1\}v}. Every other coefficient vector with value v differs from it by ker B, which is orthogonal to im B}. Its squared norm therefore equals that minimum plus the squared norm of the difference. It follows that \[
G_{N,y}^\sigma=(C_{N,y}^\sigma)^{-1},\qquad
\boxed{H_{K,N}^\sigma=I_{K,y}^*(C_{N,y}^\sigma)^{-1}I_{K,y}.}
\tag{KG19}
\] These are JM26--JM27 with the coordinate transport specified. The inverse is taken on the complete q-dimensional value space before restriction to the actual m\_K-dimensional kernel. In general \texttt{I\_K\^{}*C\^{}\{-1\}I\_K} is different from \texttt{(I\_K\^{}*C\ I\_K)\^{}\{-1\}}; only the former is proved and used here. Every column through degree N is present in KG18, so every earlier scalar relation enters the minimum.

\subsection{5. The actual change to the joint kernel metric}\label{the-actual-change-to-the-joint-kernel-metric}

Keep the same section at all four cutoffs, as in J3--J11. The already proved joint form enclosure and nesting give \[
0\le\mathcal R\log\det H_{K,N}^J\le2m_KW,
\qquad W=k\log(\tau_\eta\Lambda_{h,2q}).
\tag{KG20}
\] By definition \texttt{Delta\_\{K,N\}=log\ det\ H\_\{K,N\}-log\ det\ H\_\{K,N\}\^{}J}. Combine KG15 with both endpoints of KG20 to obtain \[
\boxed{\begin{aligned}
\mathcal R\log\det H_{K,N}^\sigma-2m_K(L_k^{\rm form}+W)
&\le\mathcal R\Delta_{K,N}\\
&\le\mathcal R\log\det H_{K,N}^\sigma+2m_KL_k^{\rm form}.
\end{aligned}}
\tag{KG21}
\] This is JM28, with its signs and asymmetric error terms unchanged. At the declared polynomial tolerances, \texttt{W=O\_h(k\ log\^{}2(q+2))}; thus both errors are o\_h(kq) by KG16. More specifically, for the actual simple one-factor divisor, \texttt{d\_max=0}, \texttt{tau\_eta=L\_h}, and W is independent of eta altogether. The polynomial-tolerance statement is therefore valid, and on this simple domain the same error bound holds for every positive eta at most one when the section is fixed across cutoffs.

The derived expression KG19 is the full finite Gamma receiver needed for the next kernel calculation. The reduction preserves a possible kq coefficient of its four-return, but does not assign that coefficient without evaluating the determinant asymptotics in the actual frame.

\subsection{Reading and use record}\label{reading-and-use-record}

The primary task source actually read was the displayed section of \path{JOINT_MINIMUM_RECEIVERS.tex} containing JM21--JM28. The prior independent proof \path{COERCIVITY_INDEPENDENT.md} supplies P2 and the exact Gamma basis and multiplier estimates; its original DLMF TeX files and precise reading record remain in \texttt{coercivity\_sources/} and its Section 11. The convolution identity and every new error estimate were derived above from the exact Gamma Euler integrals and those already checked inputs. No additional literature, numerical experiment, rendering, build, or publication is represented as performed.

\end{document}
